\begin{lemma}[Two-color degree bound]
If $\varphi$ is a proper partial edge-coloring, then every node of
$I_G(\varphi;\alpha,\beta)$ has degree at most two.
\end{lemma}
\begin{proof}
Use the incidence graph from \noderef{TwoColorIncidenceGraph}.  First let
$z$ be a vertex-node $v$.  Its neighbors are precisely the identities of
edges incident with $v$ which have color $\alpha$ or $\beta$.  There is at
most one incident edge of color $\alpha$: two distinct such edges would be
adjacent in the sense of \noderef{AdjacentEdges} and would contradict
properness as defined in \noderef{IsProperEdgeColoring}.  The same argument
gives at most one incident edge of color $\beta$.  Their union therefore
has at most two elements (and at most one if $\alpha=\beta$).

Now let $z$ be an edge-node $e$.  If $e$ has neither color, it is isolated
in the incidence graph.  Otherwise its neighbors are precisely the
vertex-nodes corresponding to entries of the unordered endpoint pair
$G.\operatorname{ends}(e)$.  An unordered pair has at most two entries, so
there are at most two such neighbors.  These two cases exhaust the disjoint
union $V\sqcup E$ and prove the degree bound.
\end{proof}
